\section{问题重述}

回焊焊接时，电路板随传送带依次经过 11 个小温区。炉内空气温度在生产前已达到稳定，而焊接区域中心温度因吸热和散热存在滞后。题目给出一组实验炉温曲线，希望据此建立可用于不同工艺参数的温度预测模型，并进一步完成速度与温区设定的优化。

11 个小温区各长 $30.5\,$cm，相邻温区之间有 $5\,$cm 间隙，炉前和炉后区域各长 $25\,$cm。车间温度为 $25\,\dC$。温区 1--5、8--9 分别联动，温区 10--11 固定为 $25\,\dC$；其余设定温度可在实验值的 $\pm10\,\dC$ 范围内调节，传送速度可在 $65\sim100\,$cm/min 内调节。炉温曲线还须满足表 \ref{tab:limits} 的工艺要求。

\begin{table}[htbp]
  \centering
  \caption{炉温曲线的工艺约束}
  \label{tab:limits}
  \small
  \begin{tabular}{lcc}
    \toprule
    指标 & 允许范围 & 单位 \\
    \midrule
    升温段温度变化率 & $0\sim3$ & \dC/s \\
    降温段温度变化率 & $-3\sim0$ & \dC/s \\
    升温时处于 $150\sim190\,\dC$ 的时间 & $60\sim120$ & s \\
    中心温度高于 $217\,\dC$ 的时间 & $40\sim90$ & s \\
    峰值温度 & $240\sim250$ & \dC \\
    \bottomrule
  \end{tabular}
\end{table}

需要完成四项任务：第一，建立焊接区域中心温度随时间变化的模型，并计算指定工况下的炉温曲线和四个位置温度；第二，在温区设定固定时求最大允许速度；第三，在温区温度和速度均可调时，使上穿 $217\,\dC$ 到峰值之间的面积最小；第四，在第三问基础上兼顾峰值两侧曲线的对称性。
